\documentclass[12pt]{article}
\usepackage{graphicx} % Required for inserting images
\usepackage[margin=0.5in]{geometry}
\usepackage{amsmath, amssymb}
\usepackage{mathrsfs}


\usepackage{soul}


% Dr. David Brown Commands
\newcommand{\ds}{\displaystyle}
\newcommand{\R}{\mathbb{R}}
\newcommand{\N}{\mathbb{N}}
\newcommand{\Q}{\mathbb{Q}}
\newcommand{\C}{\mathbb{C}}


% Commands to get script r
\usepackage{calligra}
\DeclareMathAlphabet{\mathcalligra}{T1}{calligra}{m}{n}
\DeclareFontShape{T1}{calligra}{m}{n}{<->s*[2.2]callig15}{}
\newcommand{\scripty}[1]{\ensuremath{\mathcalligra{#1}}}

% Unit commands
\newcommand{\degree}{\text{°}}
\newcommand{\unitSpeed}{\frac{\text{m}}{\text{s}}}
\newcommand{\unitAcceleration}{\frac{\text{m}}{\text{s}^2}}

% Constant commands
\newcommand{\avogadro}{6.022\times10^{23}}

\newcommand{\boltzmannConstK}{1.381\times10^{-23}\frac{\text{J}}{\text{K}}}
\newcommand{\boltzmannConstInEV}{8.617\times10^{-5}\frac{\text{eV}}{\text{K}}}
\newcommand{\planckConst}{6.626\times10^{-34}\text{J}\cdot\text{s}}

\newcommand{\atmP}{1.01325\times10^5\,\text{Pa}}

% Commands to easily write partial derivatives
\newcommand{\partDeriv}[2]{\frac{\partial #1}{\partial #2}}
\newcommand{\doublePartDeriv}[2]{\frac{\partial^2 #1}{\partial^2 #2}}


\title{Problem 6.49}
\author{Gabriel Decker}


\begin{document}



\maketitle
\begin{flushleft}


For this problem, we will compute the following for a mole ($n=1\,\text{mol}$) of Nitrogen gas at room temp ($T=300\,\text{K}$) and 1 atm of pressure ($P=\atmP$):
Energy, Enthalpy, Helmholtz free energy, Gibbs free energy, entropy, and chemical potential.
We also know that $Z_e=1$, $\epsilon=0.00025\,\text{eV}$, and the mass of a Nitrogen gas molecule is $m=4.65186\times10^{-26}$.
The relevant equations for these properties are the following.


\begin{equation}
    U=U_\text{int}+\frac{3}{2}nRT
\end{equation}
\begin{equation}
    H=U+PV=U+nRT
\end{equation}
\begin{equation}
    G=U-TS+PV=U-TS+nRT
\end{equation}
\begin{equation}
    F=-nRT[\ln(\frac{V}{v_Q})-\ln(N)+1]+F_\text{int}
\end{equation}
\begin{equation}
    S=nR\left[\ln\left(\frac{VZ_eZ_\text{rot}}{Nv_Q}\right)+\frac{7}{2}\right]
\end{equation}
where
\begin{equation}
    F_\text{int}=-nRT\ln(Z_eZ_\text{rot})
\end{equation}
and
\begin{equation}
    v_Q=\left(\frac{h}{\sqrt{2\pi mkT}}\right)^3
\end{equation}
\begin{equation}
    \mu=-kT\ln\left(\frac{VZ_eZ_\text{rot}}{Nv_Q}\right)
\end{equation}

These come from section 6.6, 6.7, problem 6.48, and previous chapters.
From the section about rotation of diatomic molecules, we know that $Z_\text{rot}=\frac{kT}{2\epsilon}$ for molecules with identical atoms. We also know that $U_\text{rot}=nRT$.\\~\\

Before solving anything, let's tackle some of the prerequisite parts.
$$Z_\text{rot}=\frac{kT}{2\epsilon}$$
$$\implies Z_\text{rot}=\frac{(\boltzmannConstInEV)\cdot(300\,\text{K})}{2\cdot(0.00025\,\text{eV})}$$
$$\implies Z_\text{rot}=51.702$$
$$v_Q=\left(\frac{h}{\sqrt{2\pi mkT}}\right)^3$$
$$\implies v_Q=\left(\frac{\planckConst}{\sqrt{2\pi(4.65186\times10^{-26}\,\text{kg})(\boltzmannConstK)(300\,\text{K})}}\right)^3$$
$$\implies v_Q=6.90354\times10^{-33}\,\text{m}^3$$
$$V=\frac{nRT}{P}$$
$$\implies V=\frac{(1\,\text{mol})(8.314\,\text{J/(mol K)})(300\,\text{K})}{\atmP}$$
$$\implies V=0.02462\,\text{m}^3$$
$$\ln\left(\frac{VZ_eZ_\text{rot}}{Nv_Q}\right)=\ln\left(\frac{(0.02462\,\text{m}^3)(1)(51.702)}{(\avogadro)(6.90354\times10^{-33}\,\text{m}^3)}\right)$$
$$\implies\ln\left(\frac{VZ_eZ_\text{rot}}{Nv_Q}\right)=19.540$$\\~\\

With that out of the way, we can finally determine the various quantities.
$$U=U_\text{int}+\frac{3}{2}nRT$$
$$\implies U=nRT+\frac{3}{2}nRT$$
$$\implies U=\frac{5}{2}nRT$$
$$\implies U=\frac{5}{2}(1\,\text{mol})(8.314\,\text{J/(mol K)})(300\,\text{K})$$
$$\implies U=6235.5\,\text{J}$$\\~\\

$$S=nR\left[\ln\left(\frac{VZ_eZ_\text{rot}}{Nv_Q}\right)+\frac{7}{2}\right]$$
$$\implies S=(1\,\text{mol})(8.314\,\text{J/(mol K)})\left[19.540+\frac{7}{2}\right]$$
$$\implies S=191.55\,\text{J/K}$$\\~\\

$$F=-nRT[\ln(\frac{V}{v_Q})-\ln(N)+1]+F_\text{int}$$
$$\implies F=-nRT[\ln(\frac{V}{v_Q})-\ln(N)+1]-nRT\ln(Z_eZ_\text{rot})$$
$$\implies F=-nRT[\ln\left(\frac{VZ_eZ_\text{rot}}{Nv_Q}\right)+1]$$
$$\implies F=-(1\,\text{mol})(8.314\,\text{J/(mol K)})(300\,\text{K})[19.540+1]$$
$$\implies F=-51231\,\text{J}$$\\~\\

$$\mu=-kT\ln\left(\frac{VZ_eZ_\text{rot}}{Nv_Q}\right)$$
$$\implies\mu=-(\boltzmannConstInEV)(300\,\text{K})(19.540)$$
$$\implies\mu=-0.505\,\text{eV}$$\\~\\

$$H=U+nRT$$
$$\implies H=6235.5\,\text{J}+(1\,\text{mol})(8.314\,\text{J/(mol K)})(300\,\text{K})$$
$$\implies H=8729.7\,\text{J}$$\\~\\

$$G=U-TS+nRT$$
$$\implies G=6235.5\,\text{J}-(300\,\text{K})(191.55\,\text{J/K})+(1\,\text{mol})(8.314\,\text{J/(mol K)})(300\,\text{K})$$
$$\implies G=-48735\,\text{J}$$



\end{flushleft}

\end{document}
